\section*{Chapter \ref{ch:b2:structure-determination} --- Structure Determination: \texorpdfstring{$^{13}$C}{13C} NMR, MS and IR Together}

\begin{solution}{exo:b2:structure-determination:1}
1,2-Dimethylbenzene: 4 (\ce{CH3}; C1/C2; C3/C6; C4/C5). 1,3-: 5 (\ce{CH3}; C1/C3; C2; C4/C6; C5).
1,4-: 3 (\ce{CH3}; C1/C4; C2/C3/C5/C6).
\end{solution}

\begin{solution}{exo:b2:structure-determination:2}
\ce{CH3CH(OH)CH2CH3}. DEPT-135: C1 (\ce{CH3}) up, C2 (CH) up, C3 (\ce{CH2}) down, C4 (\ce{CH3}) up.
DEPT-90: C2 only.
\end{solution}

\begin{solution}{exo:b2:structure-determination:3}
209: ketone carbonyl; 170: ester (or acid) carbonyl; 129: \omterm{def:b2:huckel:aromatic}{aromatic} carbon; 64: carbon bonded to an
oxygen; 14: alkyl \ce{CH3}.
\end{solution}

\begin{solution}{exo:b2:structure-determination:4}
$(2 \times 8 + 2 - 8)/2 = 5$: a benzene ring (4) and a \ce{C=O} (1). For instance methyl benzoate,
\ce{C6H5COOCH3}, or 4-methylbenzoic acid.
\end{solution}

\begin{solution}{exo:b2:structure-determination:5}
One unsaturation, a ketone (\qty{1715}{cm^{-1}}, line at 209, quaternary); one \ce{CH2} and two
\ce{CH3}, four different carbons: butan-2-one, \ce{CH3COCH2CH3}.
\end{solution}

\begin{solution}{exo:b2:structure-determination:6}
Carbon: pentan-2-one shows five lines, symmetric pentan-3-one three. Protons: pentan-2-one has a
three-proton singlet (\ce{CH3CO}); pentan-3-one only a quartet and a triplet. \omterm{def:b2:mass-spec-atomic:spectrum}{Mass spectrum}: pentan-2-one
has its \omterm{def:b2:mass-spec-atomic:spectrum}{base peak} at 43 (\ce{CH3CO+}) and a McLafferty ion at 58; pentan-3-one has its \omterm{def:b2:mass-spec-atomic:spectrum}{base peak} at 57
(\ce{C2H5CO+}).
\end{solution}

\begin{solution}{exo:b2:structure-determination:7}
\ce{C4H8O2}, 88: the quartet--triplet pair is an \ce{OCH2CH3} (the \ce{CH2} at 4.12 is on oxygen), the
singlet a \ce{CH3CO}. Ethyl ethanoate, \ce{CH3COOCH2CH3}.
\end{solution}

\begin{solution}{exo:b2:structure-determination:8}
1,2-Dichlorobenzene: three lines (C1/C2, C3/C6, C4/C5). 1,4-Dichlorobenzene: two (C1/C4, the four
CH). DEPT-90 shows two CH lines for the first, one for the second.
\end{solution}

\begin{solution}{exo:b2:structure-determination:9}
C2 is a stereocentre. Replacing one or the other proton of C3 by a deuterium gives two
diastereoisomers, not enantiomers: the two protons are diastereotopic, in different surroundings.
They can have different shifts and couple to each other as well as to their neighbours, so the
\ce{CH2} gives a complex multiplet instead of a simple quintet.
\end{solution}

\begin{solution}{exo:b2:structure-determination:10}
Five unsaturations; a monosubstituted ring (three CH lines and one C line at 137.0, five \omterm{def:b2:huckel:aromatic}{aromatic}
H); a ketone (200.8, C); an ethyl group (\ce{CH2} 31.8 quartet at 2.98 next to the carbonyl, \ce{CH3}
8.3). 1-Phenylpropan-1-one, \ce{C6H5COCH2CH3}. Conjugation with the ring delocalises the \ce{C=O}
$\pi$ electrons and weakens the bond: the band falls below the \qty{1715}{cm^{-1}} of a saturated ketone.
% ledger: ir:CO-ketone
\end{solution}

\begin{solution}{exo:b2:structure-determination:11}
Pentanoic acid: the very broad \ce{O-H} band from 2500 to \qty{3300}{cm^{-1}} and a proton near 11--12
ppm. Methyl butanoate: a three-proton singlet in the range of protons on a carbon bonded to oxygen
(3.3--4.5). Ethyl propanoate: two quartets, one on oxygen and one next to the carbonyl. Propyl
ethanoate: a three-proton singlet next to the carbonyl (2.0--2.4) and a triplet on oxygen.
% ledger: ir:OH-acid, nmrr:acid, nmrr:alcohol, nmrr:ketone
\end{solution}

\begin{solution}{exo:b2:structure-determination:12}
Both lines vanish in \omterm{def:b2:structure-determination:dept}{DEPT}, so both are quaternary, and \omterm{def:b2:structure-determination:dept}{DEPT} alone cannot tell them apart. The shifts
can: ester carbonyls lie near 167--171 and \omterm{def:b2:huckel:aromatic}{aromatic} carbons between 128 and 137 in the chart of this
chapter. A ring carbon at 166.5 and a carbonyl at 130.6 would both lie far outside their classes; the
correct assignment is the reverse.
\end{solution}

\begin{solution}{pb:b2:structure-determination:1}
\textbf{1.} $3.1/31.7 \approx 9.8~\%$; $9.8/1.11 \approx 9$ carbons.
\textbf{2.} Even mass: no nitrogen (or two).
\textbf{3.} $150 - 108 = 42$: \ce{H10O2} fits (10 + 32): \ce{C9H10O2}.
\textbf{4.} Ten carbons would give $M+1 \approx 11~\%$ of $M$, not 9.8~\%; and the IR shows a
carbonyl, which needs a second oxygen besides the \ce{C-O} band.
\textbf{5.} $(2 \times 9 + 2 - 10)/2 = 5$.
\textbf{6.} $108 + 10 \times 1.007825 + 2 \times 15.994915 \approx 150.0681$.
\textbf{7.} A \ce{C=O} stretch at the position of a saturated ester (near \qty{1735}{cm^{-1}}).
\textbf{8.} The \ce{C-O} stretch of the ester.
\textbf{9.} \omterm{def:b2:huckel:aromatic}{Aromatic} (or alkene) \ce{C-H} bonds.
\textbf{10.} An \ce{O-H} group: no alcohol, no carboxylic acid.
\textbf{11.} No: a carbonyl conjugated with a ring absorbs at lower wavenumber; here it sits where
saturated esters do, so a non-conjugating group separates it from the ring.
\textbf{12.} 170.70: ester \ce{C=O}; 136.14: ring carbon bearing the substituent; 128.56 and 128.24:
\omterm{def:b2:huckel:aromatic}{aromatic} CH; 66.24: \ce{CH2} bonded to oxygen; 20.82: \ce{CH3}.
\textbf{13.} 66.24, the \ce{CH2}.
\textbf{14.} 170.70 and 136.14, the two carbons without hydrogen.
\textbf{15.} 128.56 and 128.24.
\textbf{16.} Four: the carbon bearing the substituent, the two ortho, the two meta and the para CH. Three
CH kinds are expected; two of them have shifts too close to be resolved at this field, and share a line.
\textbf{17.} Not necessarily: heights are not counts (\cref{prop:b2:structure-determination:intensities}),
though here the line may well contain two overlapping CH kinds.
\textbf{18.} 7.33: the five \omterm{def:b2:huckel:aromatic}{aromatic} protons; 5.09: the \ce{OCH2}; 2.06: the \ce{CH3CO}.
\textbf{19.} Their neighbouring atoms (the ring carbon and the oxygen for the \ce{CH2}, the carbonyl
carbon for the \ce{CH3}) carry no hydrogen.
\textbf{20.} It is bonded to the electronegative oxygen of the ester and to the ring: both deshield it.
\textbf{21.} \ce{C6H5CH2OCOCH3}, benzyl ethanoate.
\textbf{22.} In the isomer the \ce{CH2} is not bonded to oxygen, so it could not appear at 5.09, and the
methyl, on oxygen, would appear in the 3.3--4.5 range, not at 2.06 next to a carbonyl.
% ledger: nmrr:alcohol, nmrr:ketone
\textbf{23.} 108: loss of 42 (ketene, \ce{CH2=C=O}) by rearrangement, the radical cation of
\ce{C7H8O}; 91: \ce{C7H7+}, the benzyl cation; 43: \ce{CH3CO+}.
\textbf{24.} Seven kinds of carbon (\ce{C=O}, the substituted ring carbon, ortho, meta and para CH,
\ce{CH2}, \ce{CH3}): \textbf{seven lines, one of them (the \ce{CH2}) pointing down} in DEPT-135.
\end{solution}
